% =====================================================================
%  Capítulo 1 — Introdução
% =====================================================================
\chapter{Introdução}
\label{cap:introducao}

\section{Contextualização}

A disponibilidade operacional de máquinas rotativas --- motores, bombas,
compressores, turbinas e redutores --- é um fator crítico em plantas
industriais de alta performance. A falha não planejada de um ativo não apenas
interrompe a produção, mas pode gerar danos secundários, riscos de segurança e
custos de reparo muito superiores aos de uma intervenção planejada. Para
mitigar esses riscos, a engenharia de manutenção evoluiu da manutenção
corretiva e preventiva baseada em calendário para a \emph{manutenção
preditiva}, na qual a condição do ativo é monitorada continuamente e as
intervenções são decididas com base em evidências do estado de degradação
\cite{scheffer2004,randall2011}.

Dentre as técnicas de monitoramento de condição, a análise de vibração é uma
das mais difundidas e informativas para máquinas rotativas. A vibração mecânica
é consequência direta das forças dinâmicas internas da máquina e, portanto,
carrega informação sobre o estado de componentes como rotores, mancais,
acoplamentos e rolamentos \cite{harris2002}. A análise espectral por
Transformada Rápida de Fourier (FFT) permite decompor o sinal de vibração e
identificar componentes de frequência cuja origem cinemática pode ser
associada a defeitos específicos \cite{randall2011,cooley1965}. Essa relação
entre componente espectral e fenômeno mecânico fundamenta o diagnóstico de
falhas em máquinas rotativas.

O avanço da chamada Indústria 4.0, contudo, trouxe um desafio paradoxal: a
digitalização e a sensorização massiva produzem um volume de dados brutos que,
se armazenado de forma indiscriminada, sobrecarrega as infraestruturas e
transforma um ativo analítico em passivo de armazenamento. Em sistemas de
turbomáquinas, frequências de interesse da ordem de 5~kHz exigem taxas de
amostragem de dezenas de milhares de amostras por segundo, multiplicadas por
centenas de pontos de medição. Esse cenário motivou o desenvolvimento de
estratégias de \emph{persistência seletiva}, em que a relevância técnica do
dado precede a sua gravação, e o dado bruto redundante é substituído por
representações compactas no domínio da frequência.

Paralelamente, o treinamento de analistas de vibração e de algoritmos de
inteligência artificial (IA) para diagnóstico depende da disponibilidade de
sinais de defeitos variados, em diferentes severidades e condições
operacionais. A aquisição desses sinais em ativos reais é lenta, cara e,
frequentemente, inviável: um defeito deve primeiro ocorrer, ser detectado e ser
documentado. A \emph{geração de sinais sintéticos} surge, então, como
alternativa para emular falhas de forma controlada e segura, permitindo testar
algoritmos e treinar analistas sem expor ativos críticos a riscos.

Este trabalho situa-se na interseção desses dois problemas. A partir de um
conjunto de diretrizes técnicas de engenharia de confiabilidade --- que
estabelecem critérios metrológicos de aquisição, um critério de persistência
seletiva denominado \(R^3\) e um catálogo de assinaturas espectrais de defeitos
--- foi especificado, projetado e implementado o \textbf{ARGUS}, um simulador de
defeitos em máquinas rotativas com geração de sinais vibracionais sintéticos,
análise espectral por FFT e persistência seletiva de leituras, acompanhado de
uma interface web para análise visual e registro de \emph{snapshots} de
defeito.

\section{Problema}

A análise do contexto evidencia dois problemas interligados que este trabalho
aborda:

\begin{enumerate}
  \item \textbf{Volume de dados brutos de vibração.} A aquisição contínua de
  sinais no domínio do tempo gera um volume massivo de dados. Em frequências
  de interesse de até 5~kHz, a taxa de amostragem mínima, segundo o critério de
  Nyquist, é da ordem de dezenas de milhares de amostras por segundo, o que
  torna o armazenamento indiscriminado inviável em escala industrial. É
  necessário reduzir o volume persistido sem perder a capacidade diagnóstica,
  conservando apenas a informação relevante --- os picos espectrais e as
  métricas de energia do sinal --- e descartando a redundância.
  \item \textbf{Falta de sinais de defeito para treinamento e pesquisa.}
  Analistas e algoritmos de IA precisam de exemplos de defeitos em diferentes
  estágios de severidade para aprendizado e validação. A obtenção desses
  exemplos em ativos reais é limitada por risco, indisponibilidade e custo,
  motivando a geração sintética de sinais cujas assinaturas espectrais sejam
  tecnicamente coerentes com a literatura de diagnóstico de vibração.
\end{enumerate}

O desafio de engenharia é, portanto, duplo: \emph{como gerar} sinais
vibracionais sintéticos cujas assinaturas espectrais representem defeitos
específicos de máquinas rotativas, e \emph{como persistir seletivamente} as
leituras de forma que apenas as transições de estado relevantes sejam
armazenadas, preservando a rastreabilidade para a análise de causa raiz (RCA).

\section{Justificativa}

A justificativa deste trabalho apoia-se em três pilares. O primeiro é
estratégico e econômico: a redução do volume de dados persistidos reduz o
custo de infraestrutura de armazenamento e direciona a análise para o que é
realmente crítico, aumentando o retorno sobre o investimento da manutenção
preditiva \cite{scheffer2004}. O segundo é metodológico: a capacidade de
emular falhas permite validar algoritmos de diagnóstico e treinar analistas sem
expor ativos reais, o que acelera o desenvolvimento de competências e de
modelos de IA. O terceiro é tecnológico: o trabalho materializa, em um sistema
funcional e testado, um conjunto de diretrizes de engenharia de confiabilidade
que, sem uma implementação de referência, permaneceriam apenas conceituais.

Além disso, o projeto foi conduzido seguindo uma metodologia de
especificação dirigida (Spec-Driven Development), o que gerou um acervo
documental --- especificação de requisitos, design, arquitetura, testes --- que
constitui evidência verificável do processo de engenharia adotado e base para a
reprodutibilidade dos resultados aqui apresentados.

\section{Objetivos}

\subsection{Objetivo geral}

Projetar e implementar um simulador de defeitos em máquina rotativa que gere
sinais vibracionais sintéticos com assinatura espectral parametrizável por tipo
de defeito e que aplique persistência seletiva de leituras --- por meio da
regra de ouro e do Paralelepípedo de Descarte no critério \(R^3\) --- com uma
interface web para análise espectral e registro de \emph{snapshots} de defeito
para fins de treinamento e pesquisa.

\subsection{Objetivos específicos}

\begin{enumerate}
  \item Modelar e implementar a geração de sinais sintéticos cujas componentes
  espectrais (ordens síncronas, harmônicos, sub-harmônicos, fracionários e
  frequências de falha de rolamento) reproduzam o catálogo de defeitos de
  máquinas rotativas adotado neste trabalho.
  \item Implementar o processamento espectral por FFT com extração de picos no
  espaço \(R^3\) (frequência, amplitude e fase), cálculo de RMS total, de ruído
  e de picos, valor DC e limiar configurável de detecção de picos.
  \item Implementar o motor de descarte (regra de ouro + Paralelepípedo de
  Descarte \(R^3\)) e a persistência seletiva com hierarquia
  Planta~$\rightarrow$~Área~$\rightarrow$~Máquina~$\rightarrow$~Ponto e
  armazenamento de leituras descartadas (\emph{trash}).
  \item Expor uma API REST e uma interface web com painel de sinal, gráfico de
  FFT com ordens normalizadas \(N\) e rotação em Hz, indicadores de tempo de
  processamento e taxa de descarte e funcionalidade de \emph{snapshot} de
  defeito.
  \item Validar o sistema por testes automatizados (unitários, de integração e
  de contrato) e por ensaio de ponta a ponta, verificando as assinaturas
  espectrais geradas e o fluxo de descarte.
\end{enumerate}

\section{Delimitação}

O trabalho delimita-se ao desenvolvimento de um \emph{simulador} de sinais, e
não de um sistema de aquisição real. As seguintes fronteiras são assumidas:

\begin{itemize}
  \item os sinais são gerados sinteticamente e não correspondem a medições em
  ativos reais; as amplitudes relativas dos perfis são parâmetros de síntese,
  conforme discutido no Capítulo~\ref{cap:fundamentacao};
  \item as frequências de falha de rolamento (BPFO, BPFI, BSF e FTF) são
  calculadas a partir de uma geometria padrão de rolamento;
  \item o ambiente de execução é local (Docker Compose), sem autenticação e sem
  ambiente de produção definidos, conforme decisões de projeto registradas no
  Capítulo~\ref{cap:desenvolvimento};
  \item o processamento é síncrono por requisição, sem filas ou lotes.
\end{itemize}

\section{Síntese da metodologia}

Do ponto de vista metodológico, o trabalho classifica-se como pesquisa
aplicada e desenvolvimento experimental, combinando engenharia de software,
processamento digital de sinais e engenharia de confiabilidade. O processo
adotado foi a especificação dirigida (Spec-Driven Development), no qual os
requisitos funcionais e não funcionais são formalizados antes da implementação,
o design é detalhado por tarefas pequenas e rastreáveis, e cada módulo é
validado por testes automatizados. A metodologia completa, incluindo materiais,
ambiente e critérios de aceitação, é descrita no Capítulo~\ref{cap:metodologia}.

\section{Organização do trabalho}

Além desta introdução, o trabalho está organizado da seguinte forma:

\begin{itemize}
  \item o \textbf{Capítulo~\ref{cap:fundamentacao}} apresenta a fundamentação
  teórica: vibração em máquinas rotativas, aquisição e processamento de sinais,
  FFT e análise espectral, catálogo de defeitos e o critério de persistência
  seletiva \(R^3\);
  \item o \textbf{Capítulo~\ref{cap:trabalhos}} discute trabalhos relacionados
  e posiciona a contribuição deste trabalho;
  \item o \textbf{Capítulo~\ref{cap:metodologia}} descreve materiais e métodos,
  incluindo o processo de engenharia e os critérios de validação;
  \item o \textbf{Capítulo~\ref{cap:desenvolvimento}} detalha a arquitetura e a
  implementação do sistema;
  \item o \textbf{Capítulo~\ref{cap:experimentos}} apresenta os experimentos e
  resultados;
  \item o \textbf{Capítulo~\ref{cap:discussao}} discute os resultados e as
  limitações;
  \item o \textbf{Capítulo~\ref{cap:conclusao}} conclui o trabalho e indica
  trabalhos futuros.
\end{itemize}
